\documentclass[11pt]{article}
\usepackage[margin=1.1in]{geometry}
\usepackage{amsmath,amssymb,amsthm,mathtools}
\usepackage[hidelinks]{hyperref}

\theoremstyle{plain}
\newtheorem{theorem}{Theorem}[section]
\newtheorem{proposition}[theorem]{Proposition}
\newtheorem{lemma}[theorem]{Lemma}
\newtheorem{corollary}[theorem]{Corollary}
\theoremstyle{definition}
\newtheorem{conjecture}[theorem]{Conjecture}
\newtheorem{definition}[theorem]{Definition}
\newtheorem{remark}[theorem]{Remark}
\newtheorem{question}[theorem]{Question}

\newcommand{\Z}{\mathbb{Z}}
\newcommand{\Q}{\mathbb{Q}}
\newcommand{\F}{\mathbb{F}}
\newcommand{\C}{\mathbb{C}}
\newcommand{\N}{\mathbb{N}}
\newcommand{\lat}{\tfrac{1}{N}\Z}

\title{Counting positive rational friezes over $\tfrac{1}{N}\Z$}
\author{Vico Bonfioli\\ \small\texttt{vicobonfioli@gmail.com}}
\date{}

\begin{document}
\maketitle

\begin{abstract}
Let $T(N,n)$ denote the number of positive rational frieze patterns of width $n$ all of
whose entries lie in $\tfrac1N\Z$. Conway and Coxeter showed $T(1,n)=C_{n-2}$. We prove
that for every prime $p\ge5$
\[
  T(p,5)=5+5C(p),\qquad C(p)=\#\{(a,u,v)\in\Z_{>0}^3 : auv=p+u+v\},
\]
an equality identifying the width-five count with the case $A=1$, $B=p$ of
$xyz=A(x+y)+B$, a problem posed by Ford in generalisation of a question of Conrey. The
order of $C(p)$ is open, so no closed form for $T(p,5)$ should be expected.

We also prove $T(N,4)=d(2N^2)$, give proved search bounds at widths $5$ and $6$, so that
the tabulated values are complete rather than observed to stabilise, and tabulate
$T(N,n)$ for widths $3$ to $6$ and, at small $N$, for widths $8$ and $9$. The rotation
action of $\Z/5$ on width-$5$ quiddity cycles is free, so $5\mid T(N,5)$. We prove
$T(N,5)=O_\varepsilon(N^{1+\varepsilon})$, and $T(p,5)=O_\varepsilon(p^{1/3+\varepsilon})$
at primes by importing the method of Conrey and Shah, while $T(p,5)\ge5\,d(p+1)$ refutes
the guess $T(N,5)\asymp d(N^2)\log^2N$.

The reduction behind the tables holds at every width: homogenising the frieze continuant
and applying the Desnanot--Jacobi identity collapses the lattice conditions to
$gUV=N^{n-4}(U+V)+M$ with $M\mid N^2R$ and $N^{n-4}R\mid gM$, where $R$ is a continuant
window empty at width $5$.

The count at a prime is machine-checked in Lean 4 from the definition of a frieze itself,
with periodicity derived rather than assumed; Section~\ref{sec:lean} says what else is
verified and what is not.
\end{abstract}

\section{Introduction}\label{sec:intro}

A frieze pattern of width $n$ is a finite array of rationals bordered by zeros and ones in
which every adjacent $2\times2$ minor equals $1$. Conway and Coxeter \cite{CC} classified
the patterns with positive integer entries: they correspond to triangulations of a convex
$n$-gon, so there are $C_{n-2}$ of them. Allowing denominators changes the problem. Fixing
$N$ and asking for the positive patterns with all entries in $\tfrac1N\Z$ gives a count
$T(N,n)$, finite because $\lat$ is discrete and Cuntz and Holm \cite{CH19} showed a discrete
subset admits only finitely many friezes of each height, and
computing it is the subject of this paper.

The point of this paper is that this count is an open problem in analytic number theory
wearing different clothes. At a prime,
\[
  T(p,5)=5+5C(p),\qquad C(p)=\#\{(a,u,v)\in\Z_{>0}^3 : auv=p+u+v\}
\]
(Theorem~\ref{thm:orbit}). The width-5 reduction is $guv=N(u+v)+M$, which is
$xyz=A(x+y)+B$ with $A=N$ and $B=M$, posed by Ford in generalisation of a question of
Conrey and recorded in \cite{Huang}; the cubic $auv=p+u+v$ above is its case $A=1$, $B=p$.
Ford appears to have published nothing further on it. The fully symmetric variant
$n=xyz+x+y+z$, also asked by Conrey, is Problem~25 of the American Institute of Mathematics
problem list of May 2026 and is the equation of Conrey and Shah \cite{ConreyShah}. The
order of $C(p)$ is not known. Counting friezes over $\tfrac1p\Z$ and estimating $C(p)$ are
the same problem.

Friezes have been related to equations of this shape before, and the difference matters.
Zhang \cite{Zhang} relates friezes to $xyz=G(x,y)$ in a different setting and direction:
positive \emph{integral} friezes of Dynkin type, where the Fontaine--Plamondon counts are
determined exactly, and where the Diophantine conclusion is the existence of infinitely many
positive solutions rather than a bound on their number. The equations here are the rational
ones over $\lat$, and the question is the counting bound asked by Conrey and, in generalised form, by Ford. Conley and
Ovsienko \cite{CO21} enumerate the positive \emph{integer} solutions of the same matrix
equation, graded by a parameter measuring the departure of a dissection from a triangulation;
that parameter and the denominator $N$ are independent, and none of the sequences below
occurs in their tables.

Theorem~\ref{thm:orbit} is an equality, not an estimate, so an asymptotic for either side
transfers to the other. Neither side is settled here. The count is
$O_\varepsilon(p^{1/3+\varepsilon})$ by the method of Conrey and Shah \cite{ConreyShah}
(Theorem~\ref{thm:primecube}), against a conjectured $p^{o(1)}$
(Conjecture~\ref{conj:order}). The lower bounds of Propositions~\ref{prop:lower}
and~\ref{prop:twofam}, $T(p,5)\ge5\,d(p+1)$ and $C(p)\ge d(p+1)+d(2p+1)$, are the $a=1$ and
$a=2$ slices; Conrey and Shah record the same two for their equation in
\cite[\S2.1]{ConreyShah}. Proposition~\ref{prop:modhyp} puts the count on the modular
hyperbola $de\equiv1\bmod p$, where the Weil error is of size $p^{1/2}$ against a quantity
of size $p^{1/3}$, which is why the standard route does not reach even the trivial bound.

For frieze theory this explains an absence. $T(N,4)=d(2N^2)$ has a closed form, counts of
friezes over $\Z/n\Z$ are multiplicative with explicit formulas, and no such formula has
been found at width $5$ over $\lat$: finding one would settle a question of Conrey.
Karpenkov, Short, van Son and Zabolotskii \cite{KSSZ} classified exactly these friezes and
left the count open. In the other direction it gives $C(p)$ a combinatorial reading, as the
count of a set of friezes, which is not how the quantity in Ford's problem arises
elsewhere.

Section~\ref{sec:data} tabulates $T(N,n)$ for widths $3$ to $6$ and for widths $8$ and $9$
at small $N$. At width $4$ the count has a closed form,
$T(N,4)=d(2N^2)$ (Theorem~\ref{thm:w4}). The reductions and search bounds behind the tables
at the larger widths are routine and are carried out in the accompanying library rather than
here.

Section~\ref{sec:lean} says which results are machine-checked and which rest on a written
proof; the per-statement map to Lean declarations is maintained with the library. Tables and computed constants are regenerated by scripts in \texttt{code/}, named at the
point of use; \texttt{code/paper\_numbers.py} collects those that have no script of their
own.

\section{Setting}\label{sec:setting}

We use the conventions of Karpenkov, Short, van Son and Zabolotskii \cite{KSSZ}.

\begin{definition}\label{def:frieze}
A \emph{rational frieze of width $n$} is a map
$F\colon\{(i,j): 0 \le i-j \le n\}\to\Q$, written $m_{i,j}$, with $m_{i,j}=0$ when
$i-j\in\{0,n\}$, with $m_{i,j}=1$ when $i-j\in\{1,n-1\}$, and satisfying the diamond rule
\[
  m_{i,j}\,m_{i+1,j+1}-m_{i,j+1}\,m_{i+1,j}=1 .
\]
It is \emph{positive} if every entry with $0<i-j<n$ is positive. The \emph{quiddity cycle}
is the row $i-j=2$, a sequence $(a_0,\dots,a_{n-1})$ read cyclically.
\end{definition}

The quiddity cycle determines the frieze: writing $r$ for the row index,
\begin{equation}\label{eq:rec}
  m_{r,j}=\frac{m_{r-1,j}\,m_{r-1,j+1}-1}{m_{r-2,j+1}},
\end{equation}
with the closing conditions that row $n-1$ is constant $1$ and row $n$ is constant $0$.

\begin{definition}\label{def:count}
For $N\ge 1$ let $T(N,n)$ be the number of positive rational friezes of width $n$ all of
whose entries lie in $\lat$. Quiddity cycles are counted as sequences, not up to rotation.
\end{definition}

Conway and Coxeter \cite{CC} showed that the positive \emph{integer} friezes of width $n$
correspond to triangulations of a convex $n$-gon, so
\begin{equation}\label{eq:catalan}
  T(1,n)=C_{n-2}=\frac{1}{n-1}\binom{2n-4}{n-2}.
\end{equation}
Finiteness of $T(N,n)$ is Corollary~3.8 of Cuntz and Holm \cite{CH19}, which gives it for any
discrete $R\subseteq\C$. Theorem~B of \cite{KSSZ} classifies the positive rational friezes
over $\lat$ in terms of minimal clockwise paths in a Farey-type graph. It classifies exactly
the objects counted here, and does not count them.
The present note is concerned with computing it.

\section{The continuant parameterisation}\label{sec:param}

Let $K_k$ denote the frieze continuant, $K_k=x_kK_{k-1}-K_{k-2}$ with $K_0=1$, $K_{-1}=0$.
Continuants and the continued-fraction geometry behind them are treated in \cite{K13}.

\begin{proposition}\label{prop:param}
A positive rational frieze of width $n\ge 4$ is determined by $a_0,\dots,a_{n-4}$, and
\begin{align*}
  a_{n-3}&=\frac{K_{n-4}(a_0,\dots,a_{n-5})+1}{K_{n-3}(a_0,\dots,a_{n-4})},\qquad
  a_{n-2}=K_{n-3}(a_0,\dots,a_{n-4}),\\
  a_{n-1}&=\frac{K_{n-4}(a_1,\dots,a_{n-4})+1}{K_{n-3}(a_0,\dots,a_{n-4})}.
\end{align*}
\end{proposition}

\begin{proof}
The closing conditions are $m_{n-1,j}=1$ and $m_{n,j}=0$ for every $j$, that is
$K_{n-2}(a_j,\dots,a_{j+n-3})=1$ and $K_{n-1}(a_j,\dots,a_{j+n-2})=0$. Expanding the second
at $j=0$ by the recursion $K_k=x_kK_{k-1}-K_{k-2}$ in its last argument and solving for
$a_{n-2}$ gives the middle display; expanding the first at $j=0$ and at $j=1$ and solving
for $a_{n-3}$ and $a_{n-1}$ gives the other two. The remaining $n-3$ closing conditions are
cyclic conjugates of these under the monodromy $M(a_0)\cdots M(a_{n-1})=-I$ and impose
nothing further.
\end{proof}

At $n=5$ this reads
\begin{equation}\label{eq:w5param}
  a_2=\frac{a_0+1}{D},\qquad a_3=D,\qquad a_4=\frac{a_1+1}{D},\qquad D=a_0a_1-1,
\end{equation}
and at $n=6$, with $D=K_3(a_0,a_1,a_2)=a_2(a_0a_1-1)-a_0$,
\begin{equation}\label{eq:w6param}
  a_3=\frac{a_0a_1}{D},\qquad a_4=D,\qquad a_5=\frac{a_1a_2}{D}.
\end{equation}

\begin{remark}\label{rem:glide}
A frieze of width $n$ has a glide symmetry exchanging row $r$ with row $n-r$. For $n=5$ the
only non-constant interior rows are $2$ and $3$, and they are exchanged, so a quiddity cycle in $\lat$
forces every entry into $\lat$. For $n=6$ row $3$ is fixed by the glide, and carries a
lattice condition of its own. This is the source of the extra divisibility constraints at
width $6$.
\end{remark}

\section{Width four}\label{sec:w4}

Write $d$ for the number-of-divisors function.

\begin{theorem}\label{thm:w4}
For every $N\ge 1$, $\;T(N,4)=d(2N^2)$.
\end{theorem}

\begin{proof}
Let $(a_0,a_1,a_2,a_3)$ be the quiddity cycle. By \eqref{eq:rec} the row $3$ entries are
$a_ja_{j+1}-1$, and row $3$ is constant $1$, so $a_ja_{j+1}=2$ for all $j$ modulo $4$.
Hence $a_0=a_2$, $a_1=a_3$ and $a_0a_1=2$. Conversely any $a>0$ with $a\cdot(2/a)=2$ yields
the frieze with quiddity $(a,2/a,a,2/a)$, whose remaining rows are forced and positive.

Write $a_0=p/N$ with $p$ a positive integer. Then $a_1=2N/p$, and $a_1\in\lat$ if and only
if $2N^2/p\in\Z$, that is $p\mid 2N^2$. Distinct divisors give distinct cycles, so
$T(N,4)$ is the number of positive divisors of $2N^2$.
\end{proof}

The sequence $d(2N^2)$ is \href{https://oeis.org/A361689}{OEIS A361689}.

\section{Width five}\label{sec:w5}

Put $a_0=p/N$ and $a_1=q/N$ with $p,q$ positive integers, and set $e:=pq-N^2$.

\begin{proposition}\label{prop:w5cond}
The parameters $(p,q)$ arise from a positive frieze of width $5$ over $\lat$ if and only if
\begin{equation}\label{eq:w5cond}
  e>0,\qquad N\mid e,\qquad e \mid N^2(p+N),\qquad e\mid N^2(q+N).
\end{equation}
\end{proposition}

\begin{proof}
By \eqref{eq:w5param}, $a_3=D=e/N^2$, $a_2=N(p+N)/e$ and $a_4=N(q+N)/e$. Positivity of the
frieze is equivalent to $D>0$, that is $e>0$. Membership of $a_3$ in $\lat$ is $N\mid e$,
and membership of $a_2$ and $a_4$ is the two divisibility conditions. By
Remark~\ref{rem:glide} no further condition is needed.
\end{proof}

\begin{lemma}[Search bound]\label{lem:w5bound}
Any solution of \eqref{eq:w5cond} satisfies $p\le N^3+2N^2$ and $q\le N^3+2N^2$.
\end{lemma}

\begin{proof}
Since $e>0$ and $e\mid N^2(q+N)$ with $N^2(q+N)>0$, we have $e\le N^2(q+N)$. Substituting
$e=pq-N^2$ gives $pq-N^2\le N^2q+N^3$, that is $q(p-N^2)\le N^3+N^2$. As $q\ge 1$, either
$p\le N^2$ or $p-N^2\le N^3+N^2$; in both cases $p\le N^3+2N^2$. The bound on $q$ is
symmetric.
\end{proof}

Lemma~\ref{lem:w5bound} makes the enumeration finite with an explicit range. Since
additionally $e\mid N^2(p+N)$ and $q=(e+N^2)/p$, one may loop over $p$ in the proved range
and over the divisors $e$ of $N^2(p+N)$, which is what \texttt{code/verify\_width5\_proved.py}
does. The resulting values of $T(N,5)$ for $N\le 40$ are recorded in
\texttt{data/T5.txt}. They are complete, not merely stable under an increasing cutoff.

\section{Symmetry at width five}\label{sec:sym}

The dihedral group $D_5$ acts on width-$5$ quiddity cycles by rotation and reversal.

\begin{theorem}\label{thm:free}
The rotation action of $\Z/5$ on positive rational width-$5$ quiddity cycles is free. In
particular $5\mid T(N,5)$ for every $N$.
\end{theorem}

\begin{proof}
The closing condition at width $5$ is
$(a_ja_{j+1}-1)(a_{j+1}a_{j+2}-1)=1+a_{j+1}$ for all $j$ modulo $5$. A cycle fixed by a
nontrivial rotation is constant, say $a_j=a$, whence $(a^2-1)^2=1+a$, that is
$a(a^3-2a-1)=0$. Since $a>0$ and $a^3-2a-1=(a+1)(a^2-a-1)$, the only positive root is the
golden ratio $(1+\sqrt5)/2$, which is irrational. So no positive rational cycle is fixed.
\end{proof}

\section{Data and questions}\label{sec:data}

Every value below is produced by a procedure whose search range is proved, so the counts are
complete rather than stable under an increasing cutoff.

\[
\begin{array}{c|ccccc}
 & n=3 & n=4 & n=5 & n=6 \\\hline
N=1 & 1 & 2 & 5 & 14\\
N=2 & 1 & 4 & 20 & 102\\
N=3 & 1 & 6 & 40 & 259\\
N=4 & 1 & 6 & 60 & 487
\end{array}
\]

Neither $T(N,5)$ nor $T(N,6)$ appears in the OEIS as a function of $N$, and neither is
multiplicative: $T(2,5)T(3,5)=800\ne 110=T(6,5)$.

At widths $8$ and $9$ the same reduction, applied with the glide symmetry as a pruning rule
and swept over the entry box of Cuntz and Holm \cite{CH19}, gives
\[
\begin{array}{c|cccc}
  N & 1 & 2 & 3 & 4\\ \hline
  T(N,8) & 132 & 2576 & 9980 & 28604\\
  T(N,9) & 429 & 13101 & &
\end{array}
\]
(\texttt{code/w89\_count.rs}). The values at $N=1$ are $C_6$ and $C_7$, as they must be, and
neither sequence appears in the OEIS. The sweep is complete over a proved box, so these are
counts rather than stable observations; the cost grows like $N^{3(n-3)}$, which is why the
table stops where it does. The reduction and the search bound behind it are routine and are
carried out in the accompanying library rather than here.

\begin{theorem}[Reduced form]\label{thm:t5red}
For every $N\ge1$,
\begin{multline*}
  T(N,5)=\#\bigl\{(u,v,M)\in\Z_{>0}^{3} \;:\; \gcd(u,v)=1,\ M\mid N^2,\\
  uv\mid N(u+v)+M,\ Nuv \mid M\bigl(N(u+v)+M\bigr)\bigr\}.
\end{multline*}
\end{theorem}

\begin{proof}
Multiplying $gu>N$ by $v>0$ gives $guv>Nv$, that is $Nu+M>0$, which holds for $u,M>0$; the
same argument disposes of $gv>N$. So both positivity conditions are automatic. For the
equation, $g$ is determined by $(u,v,M)$, so requiring some $g$ with $guv=N(u+v)+M$ is
requiring $uv\mid N(u+v)+M$. Given that, multiplying $N\mid gM$ by $uv$ gives
$Nuv\mid M(N(u+v)+M)$, and the two are equivalent because $uv\ne0$.
\end{proof}

\begin{proposition}[Hyperbola bound]\label{prop:hyp}
Every solution counted by Theorem~\ref{thm:t5red} satisfies $(u-N)(v-N)\le N^2+M\le 2N^2$,
and if $u\le v$ then $u<3N$.
\end{proposition}

\begin{proof}
Since $uv$ divides the positive integer $N(u+v)+M$ it is at most that integer, and
$(u-N)(v-N)=uv-N(u+v)+N^2$, so $(u-N)(v-N)\le N^2+M$; and $M\le N^2$ because $M\mid N^2$. If
$u\le N$ then $u<3N$. Otherwise $0<u-N\le v-N$, so $(u-N)^2\le(u-N)(v-N)\le 2N^2$, whereas
$u\ge 3N$ would force $(u-N)^2\ge 4N^2$.
\end{proof}

The change is one of order: the smaller parameter is bounded \emph{linearly} in $N$, where
Lemma~\ref{lem:w5bound} gives only $N^3+2N^2$ for the original parameters. The bound
$(u-N)(v-N)\le2N^2$ is attained for every $2\le N\le24$ (\texttt{code/t5\_bounds.py}).

\begin{corollary}\label{cor:t5upper}
$\displaystyle T(N,5)\;\le\;2\sum_{M\mid N^2}\ \sum_{1\le u<3N} d(Nu+M)$, and consequently
$T(N,5)=O_{\varepsilon}\bigl(N^{1+\varepsilon}\bigr)$.
\end{corollary}

\begin{proof}
By Proposition~\ref{prop:hyp} the smaller of $u,v$ is less than $3N$; fix it as $u$ and let
$v$ run over the divisors of $Nu+M$, which is legitimate because $v\mid N(u+v)+M$ and
$v\mid Nv$ together give $v\mid Nu+M$. Each unordered pair is counted at most twice. For the
second statement, $d(n)=O_\varepsilon(n^{\varepsilon})$ gives $d(N^2)=O_\varepsilon(N^{\varepsilon})$
and $d(Nu+M)=O_\varepsilon(N^{\varepsilon})$ uniformly for $u<3N$ and $M\le N^2$, so the double
sum is $O_\varepsilon(N^{1+\varepsilon})$.
\end{proof}

No matching lower bound of the same order appears to exist.

\begin{definition}\label{def:AB}
Let $p$ be prime. In the parametrisation of Theorem~\ref{thm:t5red} the parameter $M$ divides
$p^2$, so $M\in\{1,p,p^2\}$, and $M=1$ cannot occur: the condition $Nuv\mid M(N(u+v)+M)$ reads
$p\,uv\mid p(u+v)+1$ there, forcing $p\mid1$. Write $A(p)$ for the number of solutions with
$M=p$ and $B(p)$ for the number with $M=p^2$, so that $T(p,5)=A(p)+B(p)$. At $M=p$ both
divisibility conditions of Theorem~\ref{thm:t5red} reduce to $uv\mid p(u+v+1)$, so
\[
  A(p)=\#\{(u,v)\in\Z_{>0}^2:\gcd(u,v)=1,\ uv\mid p(u+v+1)\}.
\]
A quiddity cycle has $M=p$ exactly when its parameters do, so this is also the count of cycles
with $M=p$ used in Theorem~\ref{thm:orbit}.
\end{definition}

\begin{proposition}\label{prop:rigid5}
Let $p,u,v\ge1$ with $\gcd(uv,p)=1$ and $uv\mid p(u+v+1)$. Then
\[
  (u,v)\in\{(1,1),\,(1,2),\,(2,1),\,(2,3),\,(3,2)\},
\]
independently of $p$. These five pairs form a single orbit of the rotation action of $\Z/5$ of
Theorem~\ref{thm:free}; that last assertion is
not proved here and nothing below depends on it.
\end{proposition}

\begin{proof}
Since $u\mid uv\mid p(u+v+1)$ and $\gcd(u,p)=1$ we get $u\mid u+v+1$, hence $u\mid v+1$;
symmetrically $v\mid u+1$. So $u\le v+1$ and $v\le u+1$. If $u=v$ then $u\mid u+1$ gives
$u=1$. If $v=u+1$ then $u\mid v+1=u+2$ gives $u\mid2$, so $(u,v)$ is $(1,2)$ or $(2,3)$;
the case $u=v+1$ is symmetric. Coprimality of $u$ and $v$ is not needed.
\end{proof}

So every further solution at a prime requires $p\mid uv$, and $T(p,5)$ stays small: it is
$160$ at $p=541$ and $130$ at $p=421$, against $T(120,5)=4170$.

By Remark~\ref{rem:glide} the width-5 glide gives $a_{j+3}=a_ja_{j+1}-1$ for every $j$
modulo $5$, which in numerators is
\begin{equation}\label{eq:w5cyc}
  p_jp_{j+1}=p\,(p_{j+3}+p) \qquad (j\in\Z/5),
\end{equation}
so $e_j=p\,p_{j+3}$ and therefore $M_j=p\,p_{j+3}/g_j$ with
$g_j=\gcd(p_j+p,\,p_{j+1}+p)$. In particular $M_j=p^2$ forces $p\mid p_{j+3}$, so the
$M$-profile is controlled by which entries of the quiddity are integers.

\begin{proposition}\label{prop:intcount}
Let $p$ be prime. The number of integer entries in the quiddity cycle of a positive
width-5 frieze over $\tfrac1p\Z$ is either $3$ or $5$, and it is $5$ exactly for the
Conway--Coxeter frieze $(2,2,1,3,1)$ and its rotations.
\end{proposition}

\begin{proof}
Put $I=\{j: p\mid p_j\}$. If $p\nmid p_{j+3}$ then the right side of \eqref{eq:w5cyc} has
$p$-valuation exactly one, so exactly one of $p_j$, $p_{j+1}$ is divisible by $p$, and
that one has valuation exactly one. Taking $j=m+2$, this says: for every $m\notin I$,
exactly one of $m+2$, $m+3$ lies in $I$. Exhausting the $32$ subsets of $\Z/5$, the only
ones with that property have cardinality $3$ or $5$, and those of cardinality $3$ are
exactly the sets $\{i,i+1,i+3\}$. If $|I|=5$ every entry is an integer, so the frieze is a
Conway--Coxeter frieze of width $5$, and there are $C_3=5$ of those, forming one rotation
orbit.
\end{proof}

\begin{theorem}[Orbit split]\label{thm:orbit}
Let $p$ be prime and write $R=T(p,5)/5$ for the number of rotation orbits. Each orbit
contains exactly two cycles with $M=p$ and three with $M=p^2$, except the orbit of the
Conway--Coxeter frieze $(2,2,1,3,1)$, where all five have $M=p$. Hence $A(p)=2R+3$ and
$B(p)=3R-3$, so that
\[
  3A(p)-2B(p)=15 \qquad\text{and}\qquad T(p,5)=\tfrac12\bigl(5A(p)-15\bigr).
\]
\end{theorem}

\begin{proof}
By Proposition~\ref{prop:intcount}, $I=\{j:p\mid p_j\}$ has three or five elements. If it
has five the frieze is $(2,2,1,3,1)$ up to rotation, and evaluating $M_j$ at each of the
five positions gives $p$ every time.

Suppose $|I|=3$. The condition established in the proof of
Proposition~\ref{prop:intcount}, that every $m\notin I$ has exactly one of $m+2$, $m+3$ in
$I$, forces $I=\{i,i+1,i+3\}$ for some $i$. That proof also records that the divisible one of
$p_j,p_{j+1}$ has valuation exactly one whenever $p\nmid p_{j+3}$; applying this at $j=i-1$
and at $j=i+1$, where $p_{i+2}$ respectively $p_{i+4}$ is prime to $p$, gives
$v_p(p_i)=v_p(p_{i+1})=1$ exactly, while $v_p(p_{i+2})=v_p(p_{i+4})=0$ because
$i+2,i+4\notin I$. Now take the five positions in turn.

For $j=i+1$ and $j=i+4$ we have $j+3\notin I$. In each case one of the two entries
$p_j,p_{j+1}$ is prime to $p$, hence so is $g_j$, and $v_p(M_j)=1+0-0=1$, so $M_j=p$.

For $j=i+2$ and $j=i+3$ the entry $p_{i+2}$ respectively $p_{i+4}$ is prime to $p$, so
again $g_j$ is prime to $p$, while $v_p(p_{j+3})=1$; thus $v_p(M_j)=2$ and $M_j=p^2$.

For $j=i$ write $p_i=p\alpha$ and $p_{i+1}=p\beta$, so $g_i=pG$ with
$G=\gcd(\alpha+1,\beta+1)$, and \eqref{eq:w5cyc} at $j=i$ gives $p_{i+3}=p(\alpha\beta-1)$
with $\alpha\beta-1>0$. If $M_i=p$ then $g_i=p_{i+3}$, that is $G=\alpha\beta-1$. Since
$G\mid\alpha+1$ this makes $\alpha\beta-1$ divide $\alpha+1$, say
$\alpha+1=(\alpha\beta-1)s$. But \eqref{eq:w5cyc} at $j=i+2$ reads
$p_{i+2}p_{i+3}=p(p_i+p)$, that is $p_{i+2}(\alpha\beta-1)=p(\alpha+1)$, so
$p_{i+2}(\alpha\beta-1)=p(\alpha\beta-1)s$, and cancelling $\alpha\beta-1\ne0$ gives
$p_{i+2}=ps$, contradicting $p\nmid p_{i+2}$. Hence $M_i=p^2$.

So two positions carry $M=p$ and three carry $M=p^2$. Counting over the $R$ orbits,
$A(p)=2(R-1)+5$ and $B(p)=3(R-1)$, and $T(p,5)=A+B=5R$.
\end{proof}

\section{The width-five count at a prime}\label{sec:prime}

No closed form for $T(p,5)$ should be expected. Setting $u=ps$ in the $M=p$ term and
writing $v+1=st$, the condition $uv\mid p(u+v+1)$ becomes $st-1\mid p+t$, and putting
$p+t=(st-1)m$ turns it into the symmetric cubic
\begin{equation}\label{eq:cubic}
  stm-t-m=p, \qquad\text{equivalently}\qquad (st-1)(sm-1)=sp+1 .
\end{equation}
The correspondence is exact.

\begin{lemma}\label{lem:acount}
Let $p\ge5$ be prime and let $C(p)$ be the number of triples $(s,t,m)$ of positive integers
with $stm=p+t+m$. Then $A(p)=5+2C(p)$.
\end{lemma}

\begin{proof}
By Definition~\ref{def:AB}, $A(p)$ counts the coprime pairs $(u,v)$ with $uv\mid p(u+v+1)$, a
condition symmetric in $u$ and $v$. If $\gcd(uv,p)=1$ it reduces to
$uv\mid u+v+1$, and Proposition~\ref{prop:rigid5} lists the five pairs satisfying it.
Otherwise $p\mid uv$, and $\gcd(u,v)=1$ forces $p$ to divide exactly one of $u$ and $v$; by
the symmetry the two halves are equinumerous, so $A(p)-5=2\,\#\{(u,v): p\mid u\}$.

For the pairs with $p\mid u$ write $u=ps$. Then $uv\mid p(u+v+1)$ reads
$sv\mid ps+v+1$, whence $s\mid v+1$; put $v=st-1$ with $t\ge1$. The condition becomes
$s(st-1)\mid s(p+t)$, that is $st-1\mid p+t$; put $p+t=(st-1)m$ with $m\ge1$. Expanding,
$stm=p+t+m$, which is \eqref{eq:cubic}.

Conversely, given $(s,t,m)$ with $stm=p+t+m$, set $u=ps$ and $v=st-1$. Then
$st-1=(p+t)/m\ge1$, so $v\ge1$, and $\gcd(u,v)=1$: indeed $\gcd(s,st-1)=1$, and $p\nmid
st-1$, since $st\equiv1 \bmod p$ would give $m\equiv t+m$, hence $p\mid t$ and then
$st\equiv0$, a contradiction. The two maps are mutually inverse, so
$\#\{(u,v):p\mid u\}=C(p)$.
\end{proof}

Checked against $A(p)$ computed directly for the primes $5\le p\le53$
(\texttt{code/a\_count.py}), with no discrepancy. Equation~\eqref{eq:cubic} also bounds the
count, and does so better than Corollary~\ref{cor:t5upper}.

The exponent can be lowered to $\tfrac13$ by importing the method of Conrey and Shah
\cite{ConreyShah}, who bound the number of representations $n=xyz+x+y+z$ by
$n^{1/3}\log n(\log\log n)^4$. Their argument uses that full symmetry forces the smallest
variable below $n^{1/3}$. Our equation $suv=p+u+v$ is symmetric only in $u$ and $v$, but
the same bound holds.

\begin{lemma}\label{lem:cuberoot}
For $p\ge4$, every solution of $suv=p+u+v$ in positive integers has
$\min(s,u,v)^3\le2p$.
\end{lemma}

\begin{proof}
If $\min(s,u,v)^3>2p$ then $suv>2p$, so $p<u+v$ and hence $suv=p+u+v<2(u+v)\le4\max(u,v)$.
Cancelling the maximum leaves a product of two of the variables below $4$. But
$\min(s,u,v)^3>2p\ge8$ forces every variable to be at least $3$, so that product is at
least $9$.
\end{proof}

\begin{theorem}\label{thm:primecube}
$T(p,5)=O_{\varepsilon}\bigl(p^{1/3+\varepsilon}\bigr)$.
\end{theorem}

\begin{proof}
By Lemma~\ref{lem:cuberoot} some variable is at most $(2p)^{1/3}$. If it is $s=t$, then
$(tu-1)(tv-1)=tp+1$ determines $(u,v)$ from a divisor of $tp+1$. If it is $u=t$ or $v=t$,
then $v(su-1)=p+u$ determines the pair from a divisor of $p+t$. Hence
\[
  C(p)\;\le\sum_{1\le t\le(2p)^{1/3}}\bigl(d(tp+1)+2\,d(p+t)\bigr),
\]
and a divisor-sum bound of Shiu type \cite{Shiu} gives
$\ll p^{1/3}\log p\,(\log\log p)^{c}$. The conclusion follows from
$T(p,5)=5+5C(p)$ of Theorem~\ref{thm:orbit}.
\end{proof}

For all $123$ primes $5\le p<700$ the ratio $T(p,5)/\sqrt p$ has mean $15.2$, and $12.0$
over $p>400$ (\texttt{code/t5\_primes\_split.py}); it does not decay on that range.
Theorem~\ref{thm:primecube} forces it to decay eventually, and the paragraph above places
the onset near $p=10^7$, so the small-prime data says nothing against it.

The conjecture is not new, and it is not ours. Conrey asked whether the number of positive
integer solutions of $xyz+x+y=n$ is $O_\varepsilon(n^{\varepsilon})$, and Ford generalised
this to $xyz=A(x+y)+B$ with $O_\varepsilon(|AB|^{\varepsilon})$ solutions; both questions
are recorded in \cite{Huang}. Our width-5 equation $guv=N(u+v)+M$ is Ford's equation with
$A=N$ and $B=M$, so Conjecture~\ref{conj:order} below is a case of Ford's problem, summed
over $M\mid N^2$. Writing $R_a(n)=\#\{(x,y)\in\N^2: axy-x-y=n\}$ for the quantity studied
in \cite{Huang}, at a prime
\[
  C(p)=\sum_{a\ge1}R_a(p), \qquad R_1(p)=d(p+1),
\]
both verified for every prime $p<200$. The second identity is Proposition~\ref{prop:lower}:
our lower bound is the $a=1$ term. Huang proves an asymptotic for $\sum_{n\le N}R_a(n)$
with $a$ fixed, which is the averaged form of the inner sum.

That the frieze count is governed by a problem of Conrey and Ford is the reason
Conjecture~\ref{conj:order} is stated and not proved here.

\begin{conjecture}\label{conj:order}
$T(N,5)=N^{o(1)}$.
\end{conjecture}

The natural sharper guess, $T(N,5)\asymp d(N^2)\log^2 N$, is false, and the reason is a
lower bound that the numerics cannot see.

\begin{proposition}\label{prop:lower}
$C(p)\ge d(p+1)$, hence $T(p,5)\ge 5\,d(p+1)$, for every prime $p$. Consequently
$T(N,5)\asymp d(N^2)\log^2 N$ fails.
\end{proposition}

\begin{proof}
Taking $u=1$ in \eqref{eq:cubic}, written as $suv=p+u+v$, leaves $sv=p+1+v$, which is
solvable exactly when $v\mid p+1$, with $s=(p+1)/v+1$. Distinct divisors give distinct
solutions, so $C(p)\ge d(p+1)$, and $T(p,5)=5+5C(p)$ by Theorem~\ref{thm:orbit}.

For the consequence, let $m_r$ be the product of the first $r$ primes, so $d(m_r)=2^r$ and
$\log m_r\sim r\log r$. By Linnik's theorem there is a prime $p\equiv-1\pmod{m_r}$ with
$\log p\ll\log m_r$. Then $m_r\mid p+1$, so $d(p+1)\ge2^r$ while $\log^2p\ll(r\log r)^2$,
and $2^r/(r\log r)^2\to\infty$. Since $d(p^2)=3$ at a prime, the asserted equivalence would
force $d(p+1)=O(\log^2p)$.
\end{proof}

The crossover is far out of computational reach: $2^r$ overtakes $(r\log r)^2$ near
$r=20$, where $m_r$ already exceeds $10^{25}$. So the flat ratios observed for $N\le3000$
are what the proposition predicts, not evidence against it. This is why the equivalence
are what the proposition predicts, not evidence against it: numerical stability over
three decades is not an asymptotic.

Multiplying $auv=p+u+v$ through by $a$ gives
\begin{equation}\label{eq:shift}
  (au-1)(av-1)=ap+1,
\end{equation}
so the solutions with leading coefficient $a$ are the divisors of the shifted number
$ap+1$ lying in the class $-1 \bmod a$, and
\begin{equation}\label{eq:cdecomp}
  C(p)=\sum_{a\ge1}\#\{D\mid ap+1 \,:\, D\equiv-1 \bmod a\}.
\end{equation}
The term $a=1$ is $d(p+1)$, the same count as the $u=1$ family of Proposition~\ref{prop:lower}
though a different family. The term $a=2$ is the
one that costs nothing: $2p+1$ is odd, so every divisor of it is odd and hence of the form
$2u-1$, and no divisor is lost to the congruence. Thus $R_2(p)=d(2p+1)$ exactly, and since
the two families carry different values of $a$ they are disjoint.

\begin{proposition}\label{prop:twofam}
$C(p)\ge d(p+1)+d(2p+1)$.
\end{proposition}

\begin{proof}
Formula~\eqref{eq:shift} at $a=1$ and $a=2$; the two families are disjoint because $a$ is
determined by the solution. Formalised as \texttt{cubic\_lower\_two\_families}.
\end{proof}

\subsection{The count on the modular hyperbola}

The decomposition~\eqref{eq:cdecomp} can be run in the other direction, eliminating $a$
instead of $u$ and $v$. Setting $d=au-1$ and $e=av-1$, identity~\eqref{eq:shift} reads
$de=ap+1$, so $de\equiv1\bmod p$ and $a=(de-1)/p$, while $a\mid d+1$ and $a\mid e+1$ are
exactly the conditions defining $u$ and $v$.

\begin{proposition}\label{prop:modhyp}
$C(p)=\#\bigl\{(d,e)\in\Z_{>0}^2 \,:\, de\equiv1 \bmod p,\ \tfrac{de-1}{p}\mid\gcd(d+1,e+1)\bigr\}$.
\end{proposition}

\begin{proof}
The maps $(a,u,v)\mapsto(au-1,av-1)$ and $(d,e)\mapsto\bigl((de-1)/p,(d+1)p/(de-1),
(e+1)p/(de-1)\bigr)$ are mutually inverse by~\eqref{eq:shift}. That the second is defined on
the whole target set, which is where the divisibilities $(de-1)/p\mid d+1$ and
$(de-1)/p\mid e+1$ are used, is checked rather than argued below, so this proposition rests
on a computation at that point.
Checked against $C(p)$ for all primes $5\le p<120$ with no discrepancy
(\texttt{code/modular\_hyperbola.py}).
\end{proof}

So $C(p)$ counts points of the modular hyperbola $de\equiv1\bmod p$ cut by a divisibility
that ties the level $a=(de-1)/p$ to $\gcd(d+1,e+1)$. The same substitution turns the sum
that every upper bound passes through into a hyperbola count with no side condition,
\begin{equation}\label{eq:hyp}
  \sum_{a\le A}d(ap+1)=\#\{(d,e)\in\Z_{>0}^2 \,:\, de\equiv1\bmod p,\ 1<de\le Ap+1\},
\end{equation}
and this locates the wall. Counts of points of $de\equiv1\bmod p$ in regions are governed by
Kloosterman sums, and Weil's bound gives them with an error $O_\varepsilon(p^{1/2+\varepsilon})$.
With $A=(2p)^{1/3}$ the main term of~\eqref{eq:hyp} is of size $A\log p$, that is
$p^{1/3+o(1)}$. The error term exceeds the main term. Kloosterman technology therefore does
not reach the trivial bound $C(p)\ll_\varepsilon p^{1/3+\varepsilon}$ of
Theorem~\ref{thm:primecube}, let alone improve on it, and the improvement would have to come
from below the square root barrier rather than from a sharper application of Weil.

The substitution~\eqref{eq:shift} is an identity over any commutative ring
(\texttt{fixed\_point\_ring}), which is why the analogue over $\F_q[t]$ inherits the
structure of the problem intact rather than simplifying it.

\section{Formalization}\label{sec:lean}

The development accompanying this paper is a Lean 4 library against Mathlib. It contains no
\texttt{sorry}, no \texttt{native\_decide} and no additional axioms: every declaration cited
below depends only on \texttt{propext}, \texttt{Classical.choice} and \texttt{Quot.sound},
which continuous integration checks on each commit.

What is machine-checked is the enumeration and the reductions it rests on. In particular the
counts are taken over Definition~\ref{def:frieze} itself rather than over a convenient class:
the definition is transcribed with no periodicity assumed (periodicity is derived, not
assumed, at widths $4$ and $5$), and the resulting sets are put in bijection with the
arithmetic conditions the search enumerates. Theorem~\ref{thm:w4} and the divisibility
statements of Section~\ref{sec:sym} are verified in that sense.

So is the count at a prime, end to end. Theorem~\ref{thm:orbit} and Lemma~\ref{lem:acount}
compose in the library into a single declaration giving $T(p,5)=5+5C(p)$ for every prime
$p\ge5$, stated over Definition~\ref{def:frieze} with periodicity derived rather than
assumed, and with $C(p)$ the full solution count of $auv=p+u+v$ rather than a count inside a
search box.

What is not machine-checked is stated as such. The external inputs of Shiu and of Linnik are
cited, not proved; and the passages to $O_\varepsilon$ in
Theorem~\ref{thm:primecube} and Corollary~\ref{cor:t5upper} are written proofs.

A per-statement table naming, for each result, the Lean declaration that verifies it or
recording that none exists, is maintained with the library rather than here, since it changes
with the library and not with the paper. Tables and computed constants in this paper are
regenerated by the scripts named alongside them.

\begin{thebibliography}{9}
\bibitem{Zhang} R.~Zhang, \emph{A positive Siegel theorem: Dynkin friezes and positive
Mordell--Schinzel}, arXiv:2503.08800.
\bibitem{CO21} C.~H.~Conley and V.~Ovsienko, \emph{Quiddities of polygon dissections and the
Conway--Coxeter frieze equation}, arXiv:2107.01234.
\bibitem{CH19} M.~Cuntz and T.~Holm, \emph{Frieze patterns over integers and other subsets of
the complex numbers}, J. Comb. Algebra \textbf{3} (2019), 153--188. arXiv:1711.03724.
\bibitem{CC} J.~H. Conway and H.~S.~M. Coxeter, \emph{Triangulated polygons and frieze
patterns}, Math. Gaz. \textbf{57} (1973), 87--94 and 175--183.
\bibitem{KSSZ} O.~Karpenkov, I.~Short, M.~van Son and A.~Zabolotskii, \emph{Classifying
integer tilings and hypertilings}, arXiv:2601.21445.
\bibitem{K13} O.~Karpenkov, \emph{Geometry of Continued Fractions}, Algorithms and
Computation in Mathematics \textbf{26}, Springer, 2013.
\bibitem{ConreyShah} B.~Conrey and N.~Shah, \emph{Which numbers are not the sum plus the
product of three positive integers?}, arXiv:2112.15551.

\bibitem{Shiu} P.~Shiu, \emph{A Brun--Titchmarsh theorem for multiplicative functions},
J. Reine Angew. Math. \textbf{313} (1980), 161--170.

\bibitem{Huang} J.~Huang, \emph{A mean value theorem for the Diophantine equation
$axy-x-y=n$}, arXiv:1108.0095. The questions of Conrey and Ford quoted above are stated in
its introduction.

\end{thebibliography}

\end{document}
